% 中文摘要：标题黑体小三、段前 40pt 段后 20pt；正文宋体小四、固定行距 20pt
% 摘要单独成页，300~500 字，不要出现图表
\xuptheadingnotoc{摘\hspace{2\ccwd}要}

随着网络攻击手段的不断演进，传统基于规则匹配的异常检测方法在检出率与误报率之间难以取得平衡。本文针对校园网环境下流量数据维度高、攻击样本稀缺的问题，设计并实现了一种基于深度学习的网络异常检测系统。系统以流量统计特征为基础构建输入向量，利用深度自编码器学习正常流量的隐含分布，并通过重构误差识别异常行为；同时引入注意力机制提升模型对关键特征的敏感度。实验结果表明，所提方法在公开数据集上的准确率、召回率均优于对比方法，能够满足实际场景中的实时检测需求。

在系统设计方面，本文完成了数据采集模块、特征工程模块、模型训练模块与可视化告警模块的设计与实现，各模块之间采用松耦合架构，便于后续功能扩展。最后对系统进行了功能测试与性能测试，测试结果表明系统运行稳定，检测延迟满足设计要求。本文的研究工作可为同类网络异常检测系统的设计与实现提供参考。

\vspace{0.5em}
\noindent{\heiti\zihao{-4} 关键词：}深度学习；网络异常检测；自编码器；流量分析；特征工程
